\par\Needspace{8\baselineskip}\bigskip
\pdfbookmark[3]{Exposition, 5}{pallas-II-article5}
\begin{center}5.\end{center}
\noindent\hspace*{\parindent}It will be appropriate to insert here a lemma whose statement we have found nowhere and which can often be put to good use.
\par\textit{Let \(Q,R,Q',R'\) be four points on the surface of a sphere, and let \(\theta\) be the angle between the great circles \(QR,Q'R'\), taking the two branches that extend from their point of intersection in the directions from \(Q\) toward \(R\) and from \(Q'\) toward \(R'\). Then we have}
\[
\begin{aligned}
&\cos QQ'\mathbin{.}\cos RR'-\cos QR'\mathbin{.}\cos RQ'\\
&\qquad=\cos\theta\mathbin{.}\sin QR\mathbin{.}\sin Q'R'.
\end{aligned}
\]
\par We can replace \(\theta\) by the arc \(SS'\), where \(S\) and \(S'\) denote the poles of the great circles \(QR,Q'R'\), with the poles situated in the same way with respect to these great \SourceLeaf{453}circles. That is, the sense of motion from \(Q\) toward \(R\) about \(S\) is the same as the sense of motion from \(Q'\) toward \(R'\) about \(S'\).
\par Applying this theorem to formulas 27—32, we find at once
\[
\begin{array}{rl}
33)\quad&Bz-Cy=\cos X\Lambda\mathbin{.}r\BareRoot{mp}\\[4pt]
34)\quad&Cx-Az=\cos Y\Lambda\mathbin{.}r\BareRoot{mp}\\[4pt]
35)\quad&Ay-Bx=\cos Z\Lambda\mathbin{.}r\BareRoot{mp}.
\end{array}
\]
\par At perihelion we have \(\mathrm{d}r=0\), \(r-p=-ae(1-e)\), \(x=r\cos XP=a(1-e)\cos XP\), \(y=a(1-e)\cos YP\), \(z=a(1-e)\cos ZP\). Thus formulas 21, 22, 23 become
\[
\begin{array}{rl}
36)\quad&CG-BH=mpe\cos XP\\[4pt]
37)\quad&AH-CF=mpe\cos YP\\[4pt]
38)\quad&BF-AG=mpe\cos ZP.
\end{array}
\]
\par If we combine them with equations 30, 31, 32, our lemma gives
\[
B(BF-AG)-C(AH-CF)=m^{\frac32}p^{\frac32}e\cos X\Pi.
\]
\par Now the same quantity becomes \(=F(AA+BB+CC)-A(AF+BG+CH)=mpF\). Thus we have
\[
\begin{array}{rl}39)\quad&F=e\BareRoot{mp}\mathbin{.}\cos X\Pi\end{array}
\]
and similarly
\[
\begin{array}{rl}
40)\quad&G=e\BareRoot{mp}\mathbin{.}\cos Y\Pi\\[4pt]
41)\quad&H=e\BareRoot{mp}\mathbin{.}\cos Z\Pi.
\end{array}
\]
